%%%
%
% $Autor: Shubhankar Dumka, Gautam Ramesh, Siddhanth Sharma $
% $Date: 2025-12-31 $
% $File: ApplicationDevelopment.tex $
% $Version: 1.0 $
%
% !TeX spellcheck = en_GB
% !TeX program = pdflatex
% !TeX encoding = utf8
%
%%%

\chapter{Application Development}
\label{ch:application_development}

\section{Knowledge Discovery in Databases (KDD)}
\label{sec:kdd}

The Knowledge Discovery in Databases (KDD) process is the backbone of the application development lifecycle, transforming raw sensor data into actionable insights. For this license plate recognition (LPR) system on the Jetson Nano, the KDD process is divided into rigorous data selection and processing stages to ensure the YOLOv8 and EasyOCR models perform reliably in the airport deployment environment.

\subsection{Data Selection}
\label{subsec:data_selection}

The selection of appropriate training and testing data is critical for the model's ability to generalize to the specific conditions of an airport drop-off zone.

\paragraph{Origin: Source of Data}
The dataset comprises a hybrid collection of images to balance general robustness with domain-specific accuracy.
\begin{itemize}
    \item \textbf{Primary Source:} A curated set of 1,243 labeled images focusing on German license plates. This dataset was selected to capture the specific font (FE-Schrift) and layout (EU strip, state stickers) of German plates.
    \item \textbf{Secondary Source:} Publicly available datasets, including the CCPD (Chinese City Parking Dataset) and OpenALPR benchmarks, were used to pre-train the model on diverse environmental backgrounds.
    \item \textbf{Target Domain Data:} Self-collected footage using the Microsoft LifeCam connected to the Jetson Nano. These images capture the specific lighting (artificial terminal lights, daylight transitions) and angles (20-25\textdegree{} depression) encountered at the installation site.
\end{itemize}
This mix ensures the model learns generic plate features while adapting to the specific local constraints.

\paragraph{Features: Key Image/Plate Properties}
The model must learn to be invariant to several visual features intrinsic to the application domain:
\begin{itemize}
    \item \textbf{Geometric Variations:} License plates appear at various angles due to the fixed camera position relative to moving vehicles. The model handles rotation and perspective distortion up to approx. 30\textdegree.
    \item \textbf{Lighting Conditions:} The dataset includes samples from varying times of day, capturing glare from headlights at night and shadows during the day.
    \item \textbf{Resolution and Scale:} Vehicles are detected at distances ranging from 2 to 10 meters, resulting in plate regions varying from 50px to 200px in width.
    \item \textbf{Contrast:} High contrast between the black characters and white background is standard, but the model must robustly handle dirty or mud-splattered plates that reduce this contrast.
\end{itemize}

\paragraph{Data Types: File Formats}
The data pipeline utilizes standard formats to ensure compatibility between the collection, labeling, and training tools:
\begin{itemize}
    \item \textbf{Images (JPEG/PNG):} Raw camera frames are stored as JPEG for storage efficiency during collection (LifeCam stream) and PNG for lossless quality during the annotation phase.
    \item \textbf{Annotations (TXT):} Annotations follow the YOLO format (class\_id, x\_center, y\_center, width, height), normalized to [0,1]. This text-based format is lightweight and natively supported by the Ultralytics training pipeline.
    \item \textbf{Metadata (SQLite/CSV):} Operational metadata (timestamps, confidence scores, detected text) is stored in a local SQLite database for the application logic and exported to CSV for performance analysis.
\end{itemize}

\paragraph{Quality: Data Vetting Process}
To ensure high model performance, the dataset underwent a rigorous cleaning process:
\begin{itemize}
    \item \textbf{Filtering:} Images with severe motion blur or where the license plate was entirely unreadable to the human eye were removed.
    \item \textbf{Annotation Correction:} Bounding boxes were manually verified using tools like Roboflow/LabelImg. Boxes were tightened to exclude vehicle bumpers and lights, reducing background noise in the training signal.
    \item \textbf{Label Consistency:} A unified labeling schema was enforced, ensuring that only valid German license plates were labeled as class `0`, while foreign or non-standard plates were excluded or categorized separately to avoid confusion.
\end{itemize}

\paragraph{Quantity: Dataset Size}
The total effective dataset for the specific German license plate task consists of:
\begin{itemize}
    \item \textbf{Total Images:} Approximately 1,243 high-quality samples.
    \item \textbf{Split:} The data is partitioned into 70\% Training (870 images), 20\% Validation (249 images), and 10\% Testing (124 images).
    \item \textbf{Annotations:} Over 1,300 license plate instances (some images contain multiple vehicles).
\end{itemize}
While compact, this dataset is highly targeted, allowing the YOLOv8-Nano model to converge quickly without the overhead of millions of irrelevant samples.

\paragraph{Fairness/Bias: Bias Analysis}
An analysis of the dataset reveals potential biases:
\begin{itemize}
    \item \textbf{Geographic Bias:} The data heavily favors German formats (DIN 74069). Performance on non-EU plates (e.g., US imports) may be lower.
    \item \textbf{Temporal Bias:} A majority of self-collected images were taken during daylight hours. To mitigate this, augmentation techniques simulating lower brightness and contrast were applied.
    \item \textbf{Vehicle Type Bias:} Passenger cars dominate the dataset. Trucks or motorcycles with different mounting positions might have lower detection rates. Mitigation involves future data collection specifically targeting these vehicle classes.
\end{itemize}

\subsection{Data Processing}
\label{subsec:data_processing}

Once selected, the data undergoes a transformation pipeline to prepare it for the neural network.

\paragraph{One Database: Data Structure}
The data is organized in a standard YOLO directory structure on the Jetson Nano and development workstations:
\begin{verbatim}
/dataset
    /train
        /images
        /labels
    /val
        /images
        /labels
    /test
        /images
        /labels
\end{verbatim}
This structure allows the training scripts to automatically map image files to their corresponding text annotations. The operational database (SQLite) is separate, storing only the inference results (Event Logs) rather than the raw training data.

\paragraph{Properties: Data Transformation}
Before entering the model, data is transformed to match the network's input layer requirements:
\begin{itemize}
    \item \textbf{Resizing:} All input images are resized to $640 \times 640$ pixels (letterbox resize) to maintain aspect ratio while matching the YOLOv8 input tensor.
    \item \textbf{Normalization:} Pixel intensity values (0-255) are normalized to [0.0, 1.0] to accelerate gradient descent convergence.
    \item \textbf{Channel Format:} Images are converted to RGB order as expected by PyTorch (OpenCV captures in BGR by default).
    \item \textbf{OCR Preprocessing:} For the recognition stage, cropped plate regions undergo grayscaling and binary thresholding to enhance character edges for EasyOCR.
\end{itemize}

\paragraph{Outliers: Handling Outliers}
Outliers that could destabilize training are managed as follows:
\begin{itemize}
    \item \textbf{Extreme Aspects:} Plates seen from extreme acute angles ($>45^\circ$) are excluded from the training set as they are unlikely to yield readable text and may confuse the bounding box regression loss.
    \item \textbf{Occlusions:} Images where more than 50\% of the plate is obscured (e.g., by a tow bar) are removed.
    \item \textbf{False Positives:} Negative mining is performed by including images of the background environment (road, signage) without license plates to teach the model what \textit{not} to detect.
\end{itemize}

\paragraph{Anomalies: Anomaly Detection}
Automated scripts check for data integrity before training begins:
\begin{itemize}
    \item \textbf{Coordinate Checks:} A script validates that all normalized bounding box coordinates $(x, y, w, h)$ are strictly within the $[0, 1]$ interval.
    \item \textbf{File Integrity:} Checks are run to ensure no image files are corrupted (zero-byte files) and that every image has a corresponding label file.
    \item \textbf{Class Index:} Verification that all class indices in the label files map correctly to the defined classes (i.e., only index `0` for license plate).
\end{itemize}

\paragraph{Augmentation: Data Augmentation Techniques}
To artificially expand the training set and improve robustness against the environmental variables identified in the Features section, the following augmentations are applied during the training pipeline:
\begin{itemize}
    \item \textbf{Mosaic Augmentation:} Combines 4 training images into one, forcing the model to learn context and detect smaller objects.
    \item \textbf{Random HSV:} Adjusts Hue, Saturation, and Value to simulate different times of day and weather conditions.
    \item \textbf{Scale and Translate:} Randomly zooms ($\pm 50\%$) and shifts the image to simulate varying vehicle distances and camera framing errors.
    \item \textbf{Flip:} Horizontal flips are \textit{not} used for the OCR task, but for the detection task (YOLO), horizontal flips can be used if the visual features of a plate (shape/color) are symmetric, though care must be taken not to confuse the model if text directionality is learned implicitly. In this pipeline, geometric augmentations are conservative to preserve the plate's recognizable structure.
\end{itemize}
